%-----------------------------------------------------------------------------------------------------------------------------------------------%
%	The MIT License (MIT)
%
%	Copyright (c) 2021 Jitin Nair
%
%	Permission is hereby granted, free of charge, to any person obtaining a copy
%	of this software and associated documentation files (the "Software"), to deal
%	in the Software without restriction, including without limitation the rights
%	to use, copy, modify, merge, publish, distribute, sublicense, and/or sell
%	copies of the Software, and to permit persons to whom the Software is
%	furnished to do so, subject to the following conditions:
%	
%	THE SOFTWARE IS PROVIDED "AS IS", WITHOUT WARRANTY OF ANY KIND, EXPRESS OR
%	IMPLIED, INCLUDING BUT NOT LIMITED TO THE WARRANTIES OF MERCHANTABILITY,
%	FITNESS FOR A PARTICULAR PURPOSE AND NONINFRINGEMENT. IN NO EVENT SHALL THE
%	AUTHORS OR COPYRIGHT HOLDERS BE LIABLE FOR ANY CLAIM, DAMAGES OR OTHER
%	LIABILITY, WHETHER IN AN ACTION OF CONTRACT, TORT OR OTHERWISE, ARISING FROM,
%	OUT OF OR IN CONNECTION WITH THE SOFTWARE OR THE USE OR OTHER DEALINGS IN
%	THE SOFTWARE.
%	
%
%-----------------------------------------------------------------------------------------------------------------------------------------------%

%----------------------------------------------------------------------------------------
%	DOCUMENT DEFINITION
%----------------------------------------------------------------------------------------
\documentclass[a4paper,12pt]{article}

%----------------------------------------------------------------------------------------
%	PACKAGES
%----------------------------------------------------------------------------------------
\usepackage[T1]{fontenc}        % correct hyphenation/encoding for accented (Spanish) text
\usepackage[spanish,es-noshorthands,es-nolists]{babel} % Spanish hyphenation, without babel's list/shorthand changes
\usepackage{parskip}
\usepackage[dvipsnames]{xcolor} % 'usenames' removed (obsolete/ignored in xcolor)
\usepackage[scale=0.9]{geometry}
\usepackage{tabularx}
\usepackage{enumitem}           % used by the joblong environment (bullet lists)
\usepackage{titlesec}
\usepackage{fontawesome5}       % social icons
\usepackage[unicode, draft=false]{hyperref}

\definecolor{linkcolour}{rgb}{0,0.2,0.6}
\hypersetup{colorlinks,breaklinks,urlcolor=linkcolour,linkcolor=linkcolour}

% centered version of the 'X' column type
\newcolumntype{C}{>{\centering\arraybackslash}X}
% left-aligned (ragged-right) X: avoids stretching short rows -> no underfull warnings
\newcolumntype{L}{>{\raggedright\arraybackslash}X}

% custom \section  (inspired by http://stefano.italians.nl/archives/26)
\titleformat{\section}{\Large\scshape\raggedright}{}{0em}{}[\titlerule]
\titlespacing*{\section}{0pt}{10pt}{10pt}

% job listing environments
\newenvironment{jobshort}[2]
    {
    \begin{tabularx}{\linewidth}{@{}l X r@{}}
    \textbf{#1} & \hfill & #2 \\[3.75pt]
    \end{tabularx}
    }
    {}

\newenvironment{joblong}[2]
    {
    \begin{tabularx}{\linewidth}{@{}l X r@{}}
    \textbf{#1} & \hfill & #2 \\[3.75pt]
    \end{tabularx}
    \begin{itemize}[nosep, after=\strut, leftmargin=1em, itemsep=3pt, label=--]
    }
    {
    \end{itemize}
    }

%----------------------------------------------------------------------------------------
%	BEGIN DOCUMENT
%----------------------------------------------------------------------------------------
\begin{document}
\pagestyle{empty}

%----------------------------------------------------------------------------------------
%	TITLE
%----------------------------------------------------------------------------------------
\begin{tabularx}{\linewidth}{@{} C @{}}
{\Huge Luciano Scola} \\[7.5pt]
\href{https://github.com/llascola}{\raisebox{-0.05\height}{\faGithub}\ llascola} \ $|$ \
\href{https://linkedin.com/in/llascola}{\raisebox{-0.05\height}{\faLinkedin}\ llascola} \ $|$ \
\href{mailto:lucianoluigiscola@gmail.com}{\raisebox{-0.05\height}{\faEnvelope}\ lucianoluigiscola@gmail.com} \ $|$ \
\href{tel:+5493415622332}{\raisebox{-0.05\height}{\faMobile}\ +54 9 341 5622-332} \\
\end{tabularx}

%----------------------------------------------------------------------------------------
%	EXPERIENCE
%----------------------------------------------------------------------------------------
\section{Resumen Profesional}
Desarrollador de software con sólida formación en Ciencias de la Computación. Especializado en backend, arquitecturas distribuidas y aplicación de fundamentos matemáticos en soluciones tecnológicas. Actualmente desarrollando productos Fintech en LB Finanzas con un enfoque en escalabilidad y automatización.

Rosario, Santa Fe, Argentina.

\section{Experiencia Laboral}

\begin{joblong}{Full Stack Developer | LB Finanzas}{Enero 2025 - Presente}
\item Diseño y desarrollo de microservicios en Go aplicando Arquitectura Hexagonal y Domain-Driven Design (DDD), con flujos de trabajo asíncronos y desacoplados basados en eventos (Pub/Sub).
\item Desarrollo de servicios transaccionales para operar en mercados de EE.UU. mediante integración con Alpaca API.
\item Implementación de sistemas de reconciliación de eventos para garantizar la integridad de datos, y \textit{cronjobs} para alertas de precios en tiempo real y ejecución automática de órdenes.
\item Automatización integral de procesos de \textit{onboarding} para empresas y gestión de documentación legal, optimizando los flujos de cumplimiento normativo (Compliance).
\item Desarrollo e implementación de interfaces frontend para la integración y consumo de los servicios backend desplegados.
\end{joblong}

\vspace{5pt}

\begin{jobshort}{Ayudante de Cátedra - UNR}{Marzo 2023 - Diciembre 2024}
Ayudante de Segunda en las cátedras de Programación I, Lenguajes Formales y Computabilidad, y Lógica en la carrera de Licenciatura en Ciencias de la Computación.
\end{jobshort}

%----------------------------------------------------------------------------------------
%	PROJECTS
%----------------------------------------------------------------------------------------
\section{Proyectos}

\begin{tabularx}{\linewidth}{ @{}l r@{} }
\textbf{MemCache} & \hfill \href{https://github.com/AgusM01/MiniMemcached}{Link del Repositorio} \\[3.75pt]
\multicolumn{2}{@{}X@{}}{Tecnologías utilizadas: C, Erlang, modelo Cliente-Servidor, Concurrencia.}  \\
\end{tabularx}

\begin{tabularx}{\linewidth}{ @{}l r@{} }
\textbf{Compilador FD4} & \hfill \href{https://github.com/llascola/FD4-compiler}{Link del Repositorio} \\[3.75pt]
\multicolumn{2}{@{}X@{}}{Tecnologías utilizadas: Haskell, Mónadas, DSL, Máquinas Abstractas, Programación Funcional.}  \\
\end{tabularx}

\begin{tabularx}{\linewidth}{ @{}l r@{} }
\textbf{Compresor de Archivos} & \hfill \href{https://github.com/llascola/compresor-huffmancode}{Link del Repositorio} \\[3.75pt]
\multicolumn{2}{@{}X@{}}{Tecnologías utilizadas: C, Codificación de Huffman.}  \\
\end{tabularx}

\begin{tabularx}{\linewidth}{ @{}l r@{} }
\textbf{Corrector Ortográfico} & \hfill \href{https://github.com/llascola/spell-checker}{Link del Repositorio} \\[3.75pt]
\multicolumn{2}{@{}X@{}}{Tecnologías utilizadas: C, Tablas Hash.}  \\
\end{tabularx}

\begin{tabularx}{\linewidth}{ @{}l r@{} }
\textbf{Generador de Sopas de Letras} & \hfill \href{https://github.com/llascola/sopa-letras-backtracking}{Link del Repositorio} \\[3.75pt]
\multicolumn{2}{@{}X@{}}{Tecnologías utilizadas: Python, C.}  \\
\end{tabularx}

%----------------------------------------------------------------------------------------
%	EDUCATION
%----------------------------------------------------------------------------------------
\section{Educación}
\begin{tabularx}{\linewidth}{@{}l L@{}}
2021 - Presente & Licenciatura en Ciencias de la Computación, 75\% completada. \textbf{UNR} \\
\end{tabularx}

%----------------------------------------------------------------------------------------
%	SKILLS
%----------------------------------------------------------------------------------------
\section{Habilidades}
\begin{tabularx}{\linewidth}{@{}l L@{}}
Nivel de Inglés        & B2 \\
Matemáticas Avanzadas  & Álgebra Abstracta, Análisis Matemático, Grafos, Probabilidad y Estadística. \\
Arquitectura de Software & Microservicios, Domain-Driven Design (DDD), Arquitectura Hexagonal, Event-Driven (Pub/Sub). \\
Desarrollo Backend     & Go (Golang), Node.js, Python, C/C++, REST APIs, SQL, Integraciones Fintech. \\
Desarrollo Frontend    & JavaScript, React, React Native. \\
Automatización         & Shell scripting, Cronjobs, Reconciliación de datos, Linux/UNIX administration. \\
Inteligencia Artificial & Conocimiento teórico. Uso de herramientas IA y desarrollo de aplicaciones IA. \\
\end{tabularx}

\end{document}
